\chapter*{Kurzfassung}
\addcontentsline{toc}{chapter}{Kurzfassung}

Die selektive Anwendung verschiedener Stile des \ac{NPR} führt in komplexen 3D-Szenen zu einer
hohen stilistischen Heterogenität.
Diese Material-Entropie kollidiert fundamental mit etablierten Echtzeit-Rendering-Architekturen, die strukturell auf
homogene Beleuchtungsberechnungen optimiert sind.
Es fehlen grundlegende empirische Systeme zur Messung der genauen hardwaretechnischen Limitierungen der jeweiligen
Basis-Rendering-Paradigmen in diesem Kontext.

Die vorliegende Arbeit liefert eine holistische Systemanalyse zur Quantifizierung dieses Konflikts.
Mithilfe des entwickelten Evaluationsframeworks \ac{HyDra} wurden die Basis-Architekturen in Form von
Single-Pass-Forward, Single-Pass-Deferred (mittels monolithischer Shader-Struktur) und Multi-Pass-Deferred
(mittels Light-Volumes) isoliert.
Die anschließende Evaluation erfolgte unter kontrollierter Skalierung von Geometrie, Beleuchtung und Stil-Entropie.

Die Messungen belegen, dass für dieses Szenario keine universell überlegene Architektur existiert.
Bei hoher Lichtkomplexität kollabiert die Forward-Pipeline unter multiplikativer Rechenlast und bestimmte Formen der
selektiven \ac{NPR}-Stilisierung sind nicht anwendbar.
Der Single-Pass-Deferred-Ansatz skaliert effizient unter komplexen Beleuchtungsszenarien, verzeichnet jedoch bei
hoher Stil-Entropie massive Leistungseinbrüche durch Warp-Divergenz und Registerdruck.
Die Multi-Pass-Deferred-Architektur löst diese Probleme, sättigt jedoch bei überlappenden Lichtquellen durch
übermäßige \ac{RMW}-Zyklen die \ac{VRAM}-Bandbreite der Grafikkarte.

Aus der Synthese dieser Architektur- und Hardware-Flaschenhälse wurde ein heuristisches Entscheidungsmodell
abgeleitet.
Dieses stellt Softwareentwicklern, neben der entwickelten Evaluationssoftware selbst, ein objektives Werkzeug zur
frühzeitigen, szenenspezifischen Architekturbewertung im Kontext selektiver Stilisierung zur Verfügung.

